% FAC-EN-U0056
% Source: fac_body.tex lines 1999--2040.
% Authority: PDF physical pp. 55--56; printed pp. 250--251 (no. 58 complete).
% English source-first translation; printed grammatical slip corrected transparently and logged.

\phantomsection\label{fac:III:2:no:58}
\textbf{58. First properties of the functor $\mathcal{A}(M)$.}

\phantomsection\label{fac:III:2:no:58:prop:3}
\textsc{Proposition 3.} \textit{The functor $M\mapsto\mathcal{A}(M)$ is exact.}

Let
\phantomsection\label{fac:III:2:no:58:formula:graded-exact-sequence}
\[
M\xrightarrow{\alpha}M'\xrightarrow{\beta}M''
\]
be an exact sequence of graded $S$-modules, and let us show that the
sequence
\phantomsection\label{fac:III:2:no:58:formula:stalk-exact-sequence}
\[
M_x\xrightarrow{\alpha}M'_x\xrightarrow{\beta}M''_x
\]
is also exact. Let $m'/Q\in M'_x$ be an element of the kernel of $\beta$.
By the \hyperref[fac:III:2:no:57:def:M-U]{definition of $M''_x$}, there exists $R\in S(x)$ such that
$R\beta(m')=0$. But then there exists $m\in M$ such that
$\alpha(m)=Rm'$, and
\phantomsection\label{fac:III:2:no:58:formula:preimage}
\[
\alpha\!\left(\frac{m}{RQ}\right)=\frac{m'}{Q},
\]
as required. (Compare \hyperref[fac:II:4:lemma:1]{no.~48, Lemma~1}.)

\phantomsection\label{fac:III:2:no:58:prop:4}
\textsc{Proposition 4.} \textit{If $M$ is a graded $S$-module and $n$ is
an integer, then $\mathcal{A}(M(n))$ is canonically isomorphic to
$\mathcal{A}(M)(n)$.}

Let $i\in I$, $x\in U_i$, and $m/Q\in M(n)_x$, with
$m\in M(n)_p$, $Q\in S(x)$, and $\deg Q=p$. Set
\phantomsection\label{fac:III:2:no:58:formula:eta}
\[
\eta_{i,x}(m/Q)=\frac{m}{t_i^nQ}\in M_x,
\]
which is legitimate since $m\in M_{n+p}$ and $t_i^nQ\in S(x)$. It is
immediate that $\eta_{i,x}:M(n)_x\to M_x$ is bijective for every
$x\in U_i$, and defines an isomorphism $\eta_i$ from
$\mathcal{A}(M(n))$ onto $\mathcal{A}(M)$ over $U_i$. Moreover,
\phantomsection\label{fac:III:2:no:58:formula:eta-transition}
\[
\eta_i\circ\eta_j^{-1}=\theta_{ij}(n)
\]
over $U_i\cap U_j$. In view of the definition of the operation
$\mathcal{F}(n)$ in \hyperref[fac:III:2:no:54]{no.~54}, and
\hyperref[fac:I:1:no:4:prop:4]{Proposition~4 of no.~4}, this proves that
$\mathcal{A}(M(n))$ is isomorphic to
$\mathcal{A}(M)(n)$.

\phantomsection\label{fac:III:2:no:58:corollary}
\textsc{Corollary.} \textit{$\mathcal{A}(S(n))$ is canonically isomorphic
to $\mathcal{O}(n)$.}

Indeed, we have already said that $\mathcal{A}(S)$ is isomorphic to
$\mathcal{O}$.

(It is also directly evident that $\mathcal{A}(S(n))$ is isomorphic to
$\mathcal{O}'(n)$, since $\mathcal{O}'(n)_x$ consists precisely of the
rational fractions $P/Q$ such that
$\deg P-\deg Q=n$ and $Q\in S(x)$.)

\phantomsection\label{fac:III:2:no:58:prop:5}
\textsc{Proposition 5.} \textit{Let $M$ be a graded $S$-module
\hyperref[fac:III:2:no:56:def:TF]{satisfying condition (TF)}. Then the algebraic
sheaf $\mathcal{A}(M)$ is coherent, and $\mathcal{A}(M)=0$ if and only if
$M\in\mathcal{C}$.}

If $M\in\mathcal{C}$, then for every $m\in M$ and every $x\in X$ there
exists $Q\in S(x)$ such that $Qm=0$: it is enough to take $Q$ of
sufficiently large degree. Thus $M_x=0$, whence $\mathcal{A}(M)=0$.

Now let $M$ be a graded $S$-module
\hyperref[fac:III:2:no:56:def:TF]{satisfying condition (TF)}. There is a
finitely generated homogeneous submodule $N$ of $M$ such that
$M/N\in\mathcal{C}$. By what precedes and
\hyperref[fac:III:2:no:58:prop:3]{Proposition~3}, the map
$\mathcal{A}(N)\to\mathcal{A}(M)$ is bijective, so it is enough to prove
that $\mathcal{A}(N)$ is coherent. Since $N$ is finitely generated, there
is an exact sequence
\phantomsection\label{fac:III:2:no:58:formula:free-presentation}
\[
L^1\longrightarrow L^0\longrightarrow N\longrightarrow0,
\]
where $L^0$ and $L^1$ are finitely generated \textit{free} modules. By
\hyperref[fac:III:2:no:58:prop:3]{Proposition~3}, the sequence
\phantomsection\label{fac:III:2:no:58:formula:sheaf-presentation}
\[
\mathcal{A}(L^1)\longrightarrow\mathcal{A}(L^0)
\longrightarrow\mathcal{A}(N)\longrightarrow0
\]
is exact. By the corollary to
\hyperref[fac:III:2:no:58:prop:4]{Proposition~4}, however,
$\mathcal{A}(L^0)$ and $\mathcal{A}(L^1)$ are isomorphic to finite direct
sums of sheaves $\mathcal{O}(n_i)$, and are therefore coherent. It follows
that $\mathcal{A}(N)$ is coherent.

Finally, let $M$ be a graded $S$-module
\hyperref[fac:III:2:no:56:def:TF]{satisfying (TF)} and such that
$\mathcal{A}(M)=0$. By what precedes, we may suppose that $M$ is finitely
generated. If $m$ is a homogeneous element of $M$, let
$\mathfrak{a}_m$ be the annihilator of $m$, that is, the set of
polynomials $Q\in S$ such that $Qm=0$; clearly $\mathfrak{a}_m$ is a
homogeneous ideal. Moreover, the assumption $M_x=0$ for every $x\in X$
implies that the zero locus of $\mathfrak{a}_m$ in $K^{r+1}$ is empty or
reduced to $\{0\}$. Hilbert's Nullstellensatz then shows that every
homogeneous polynomial of sufficiently large degree belongs to
$\mathfrak{a}_m$. Applying this to a finite system of generators of $M$,
we conclude at once that $M_p=0$ for all sufficiently large $p$, which
completes the proof.

Combining \hyperref[fac:III:2:no:58:prop:3]{Propositions~3} and
\hyperref[fac:III:2:no:58:prop:5]{5}, we obtain:

\phantomsection\label{fac:III:2:no:58:prop:6}
\textsc{Proposition 6.} \textit{Let $M$ and $M'$ be two graded
$S$-modules \hyperref[fac:III:2:no:56:def:TF]{satisfying condition (TF)},
and let $\varphi:M\to M'$ be a
homomorphism. The homomorphism}
\phantomsection\label{fac:III:2:no:58:formula:A-phi}
\[
\mathcal{A}(\varphi):\mathcal{A}(M)\longrightarrow\mathcal{A}(M')
\]
\textit{is injective (respectively surjective, bijective) if and only if
$\varphi$ is $\mathcal{C}$-injective (respectively
$\mathcal{C}$-surjective, $\mathcal{C}$-bijective).}
